\documentclass[11pt]{article}
\usepackage[T1]{fontenc}
\usepackage[english]{babel}
\usepackage{lmodern,amsmath,amssymb,amsthm,booktabs,microtype}
\usepackage[margin=1in]{geometry}
\usepackage[colorlinks=true,linkcolor=black,citecolor=blue,urlcolor=blue]{hyperref}
\newtheorem{theorem}{Theorem}[section]
\newtheorem{proposition}[theorem]{Proposition}
\newtheorem{corollary}[theorem]{Corollary}
\newcommand{\NN}{\mathbb Z_{\geq 0}}
\newcommand{\rev}{\operatorname{rev}}
\title{Counting Permutations of Consecutive Triple Sums:\\[2pt]
OEIS A180206 through \(n=26\)}
\author{Carlo Corti\\Independent researcher\\\texttt{carlo.corti@outlook.com}}
\date{November 2026}
\begin{document}
\maketitle

\begin{abstract}
Which permutations of \(0,\ldots,n-1\) occur as consecutive triple
sums of nonnegative integers? The zero sum fixes three entries of the
realizing vector. Reading outwards from those entries gives two words on
complementary label sets and determines the entire vector uniquely.
We count the words by a weighted subset recurrence whose continuation state
is an ordered pair. A congruence modulo three restricts each subset to one residue class
of pairs (117 of 351 pairs at \(n=26\)). A second enumeration allows repeated labels and removes
them by inclusion--exclusion. The two computations agree for every
\(1\le n\le26\), extending the historical sixteen-term prefix of
OEIS A180206 by ten terms, with \(a(26)=1688253382074096582\).
A separate membership test based on cumulative differences checks the
definition exhaustively in Python through \(n=8\) and in a separate C
program through \(n=11\). A forward count from all possible starting pairs,
carried out after the original computation, agrees through \(n=21\).
Explicit permutations distinguish
the admissible sets for windows of lengths two, three and four.
\end{abstract}

\noindent\textbf{2020 Mathematics Subject Classification.}
Primary 05A05. Secondary 05A15, 05A19.

\smallskip
\noindent\textbf{Key words and phrases.}
Permutation, consecutive sums, nonnegative realization, subset dynamic
programming, inclusion--exclusion, OEIS A180206.

\section{The inverse problem}

Let \(p=(p_0,\ldots,p_{n-1})\) be a permutation of \(0,\ldots,n-1\),
where \(n\ge1\). A \emph{realization} of \(p\) is a vector
\(x=(x_0,\ldots,x_{n+1})\in\NN^{n+2}\) such that
\begin{equation}
 p_i=x_i+x_{i+1}+x_{i+2}\qquad(0\le i<n).
 \label{eq:def}
\end{equation}
Write \(a(n)\) for the number of permutations admitting a realization. Reversals count separately
whenever they differ.
The entries of \(x\) may repeat. For example, \((2,0,1,3)\) is realized by
\((2,0,0,0,1,2)\). The position of the zero already determines the three
central entries in this example. Their two neighbours, and then every
remaining entry, can be read off from the prescribed sums.

R.~H. Hardin introduced this sequence as OEIS A180206, with initial
index 1 and author date 15 August 2010. The historical entry used at
the start of this computation lists \(a(1),\ldots,a(16)\), ending at
\(612195224\). On 18 November 2026 the OEIS text entry, its synthesized
b-file and the official OEIS data mirror gave the same sixteen terms~\cite{hardin}. Table~\ref{tab:terms} gives
the ten subsequent values computed here. The comparison throughout is with
this documented baseline.

The constructions for adjacent pairs and for windows of four in
A180205 and A180207~\cite{hardintwo,hardinfour,two,four} suggest counting outwards
from zero. Three adjacent summands require keeping two successive entries
of the realization. The distinction matters even for words ending with
the same sum of those two entries, as shown in Section~\ref{sec:weights}.
There is also a forward problem for exactly the same linear map:
Hardin's A229561 counts its distinct output arrays when the \emph{input}
is a permutation~\cite{forward}. In \eqref{eq:def} the permutation
constraint is on the output and nonnegativity is the constraint on the input.

\section{Two words determined by the zero}\label{sec:zero}

Start with the pair \((0,0)\). Reading a positive label \(v\) in state
\((a,b)\) appends the realization entry \(c=v-a-b\) and replaces the pair
by \((b,c)\). A word is \emph{admissible} if each appended entry is
nonnegative. For a finite set \(S\) of distinct positive integers let \(F(S)\)
count the admissible orderings of precisely its elements. The empty word
contributes \(F(\varnothing)=1\). Write \(U_m=\{1,\ldots,m\}\), with
\(U_0=\varnothing\).

\begin{theorem}\label{thm:split}
Each feasible permutation has exactly one realization. For \(m=n-1\)
there is a bijection between the feasible permutations and the triples
\((S,L,R)\) where \(S\subseteq U_m\), \(L\) is an admissible ordering of
\(S\), and \(R\) an admissible ordering of \(U_m\setminus S\). Consequently
\begin{equation}
 a(n)=\sum_{S\subseteq U_m}F(S)F(U_m\setminus S).
 \label{eq:split}
\end{equation}
\end{theorem}

\begin{proof}
Let \(p_k=0\). Nonnegativity in \eqref{eq:def} forces
\(x_k=x_{k+1}=x_{k+2}=0\). To move right, use \(p_{k+1}\) and the
last two zeros to determine \(x_{k+3}\), then use \(p_{k+2}\) and
the last two known entries, and so on. To move left, start with
\(x_{k+1},x_k\) and use \(p_{k-1},p_{k-2},\ldots,p_0\) in that order.
For a left window \(p_i\) the pair is read as \((x_{i+2},x_{i+1})\),
so the new entry is \(x_i=p_i-x_{i+2}-x_{i+1}\) and the new pair is
\((x_{i+1},x_i)\). This determines every
entry of any realization and produces admissible words \(L,R\).

Conversely, starting at the zero pairs, let \(y_L,y_R\) be the lists
of entries appended by the words \(L,R\), respectively. Concatenate
\(\rev(y_L),(0,0,0),y_R\). Writing \(k=|S|\), the length is \(k+3+(m-k)=n+2\). The
\(j\)-th outward label lies in window \(k-j\) on the left
and \(k+j\) on the right, for the respective word lengths.
The windows on each side have exactly the prescribed sums by construction.
The middle window has sum zero, and no window meets both appended lists.
The resulting permutation is \((\rev(L),0,R)\), from which one recovers
\(S,L,R\). This includes an empty word on either side. At \(n=1\)
both are empty and the sole realization is \((0,0,0)\).
\end{proof}

For a fixed \(S\), choosing \(L\) and \(R\) fills every position in the
permutation. The summation already runs over the actual left label sets.
For \(m>0\) one may sum over \(S\) omitting \(m\) and double the result,
since the summands for complementary sets agree.

\begin{corollary}\label{cor:bounds}
For \(n>1\), \(a(n)\) is even. For \(n\ge1\),
\(a(n+1)\ge2a(n)\) and \(2^{n-1}\le a(n)\le n!\).
\end{corollary}
\begin{proof}
Reversing a realization reverses its windows. An ordered permutation with
two or more distinct entries differs from its reversal, giving pairs.
For the second assertion, append the label \(n\) to a feasible permutation
of length \(n\). Its last two realization entries sum to at most
\(p_{n-1}\le n-1\), so appending \(x_{n+2}=n-x_n-x_{n+1}\) gives a
nonnegative entry. Reversal gives the corresponding construction at the
left end. Each map is injective by deleting \(n\), and their images are
disjoint by the position of \(n\). Iterate from \(a(1)=1\). The upper bound
counts all permutations.
\end{proof}

\subsection{A membership test using the original windows}\label{sec:direct}

A single permutation can be tested in linear time without locating its
zero. The following criterion also applies to an arbitrary list of
nonnegative integer sums \(p_0,\ldots,p_{n-1}\).

\begin{proposition}\label{prop:direct}
For \(r=0,1,2\), form a list starting at zero and adding, in increasing
order of \(i\), the differences \(p_{i+1}-p_i\) with \(0\le i<n-1\)
and \(i\equiv r\pmod3\). Let \(m_r\) be the negative of the minimum
of the resulting partial sums, including the initial zero. A realization
exists if and only if \(\sum_r m_r\le p_0\).
\end{proposition}
\begin{proof}
Subtract consecutive equations in \eqref{eq:def} to get
\(x_{i+3}-x_i=p_{i+1}-p_i\). Every entry in residue class \(r\)
is its initial value \(x_r\) plus the corresponding partial sum.
All entries in that class are nonnegative exactly when \(x_r\ge m_r\).
For \(n=1\) the three lists contain only their initial zero.
Choose integers satisfying the three lower bounds and
\(x_0+x_1+x_2=p_0\). This is possible exactly under the stated inequality.
The first window is then correct and the differences imply the rest.
If \(p\) is a feasible permutation, the slack must be zero: positive
slack could all be added to \(x_0\) or all to \(x_1\), giving
different realizations, contrary to Theorem~\ref{thm:split}.
\end{proof}

\section{Counting tails with multiplicity}\label{sec:weights}

The use of a visited subset and terminal data follows the classical
subset dynamic programs of Bellman and Held and Karp~\cite{bellman,held}.
Here the terminal data consist of two realization entries, and a state
accumulates the number of admissible orderings.

Let \(W(S,a,b)\) count the admissible words on \(S\) whose terminal pair
is \((a,b)\). Initially only \(W(\varnothing,0,0)=1\) is nonzero.
An unused positive label \(v\ge a+b\) carries the entire weight to
\(W(S\cup\{v\},b,v-a-b)\). Grouping words by their final label gives
the equivalent destination-row recurrence, for \(S\ne\varnothing\),
\begin{equation}
 W(S,b,c)=\sum_{\substack{v\in S\\v\ge b+c}}
       W(S\setminus\{v\},v-b-c,b),\qquad
 F(S)=\sum_{b,c\ge0}W(S,b,c).
 \label{eq:pull}
\end{equation}
Indeed the last label and destination pair determine the predecessor
pair uniquely. Removing the last label strictly decreases the cardinality.
Induction from the empty word therefore counts every admissible ordering
once, including distinct histories with the same terminal pair.

For example, \((2,4,5,3)\) and \((3,4,5,2)\) both end at \((1,0)\) on
the same label set. A Boolean state loses their multiplicity.
Keeping just \(a+b\) loses information even if each state is weighted.
On \(U_5\) the words \((1,2,4,3,5)\) and \((1,2,3,5,4)\) end at
\((0,3)\) and \((3,0)\). Continuing in \(U_7\), after appending 7 the pairs become \((3,4)\)
and \((0,4)\). Only the second now admits label 6.

\begin{proposition}\label{prop:slice}
For \(S\ne\varnothing\), every reachable pair satisfies
\begin{equation}
 a+b\le\max S,\qquad 2a+b\equiv\sum_{v\in S}v\pmod3.
 \label{eq:res}
\end{equation}
Every finite set of positive labels admits its increasing word.
\end{proposition}
\begin{proof}
If the last label read is \(v_h\), nonnegativity gives \(a+b\le v_h\le\max S\). The initial
pair has residue zero, and for \(c=v-a-b\) we have
\(2b+c=(2a+b)+v-3a\). This proves the congruence by induction.
For the increasing word \(v_1<\cdots<v_h\), denote the appended entries
by \(y_1,\ldots,y_h\), with \(y_{-1}=y_0=0\). We have \(y_1=v_1\)
when \(h>0\), and for \(1\le j<h\),
\(y_{j+1}=y_{j-2}+v_{j+1}-v_j\). All appended entries are positive.
\end{proof}

The restriction \eqref{eq:res} acts on pairs within a subset. For instance,
\(S=\{3\}\) admits only \((0,3)\), although \((0,0)\) passes both
restrictions. The coefficient also depends on the values of the labels.
For a word of at most three distinct positive labels the appended entries
are its first label and its successive differences, so its only admissible
ordering is increasing. In particular \(F(\{1,2,3\})=F(\{1,2,4\})=1\).
For four labels the first three must still increase. If the word is
\((u,v,w,z)\), the fourth entry is \(u+z-w\). Choosing the last label
and testing this inequality gives
\[
 \begin{split}
 F(\{1,2,3,4\})&=2 \quad\text{from }(1,2,3,4),(1,2,4,3),\\
 F(\{1,2,3,5\})&=1 \quad\text{from }(1,2,3,5).
 \end{split}
\]
Thus \(a(5)=14+2\cdot2=18\): the fourteen nonempty proper subsets
of \(U_4\) each contribute one in \eqref{eq:split}.

\subsection{Packed rows and arithmetic}\label{sec:impl}

A job with maximum length \(N\) first evaluates \eqref{eq:pull} over
subsets of \(U_M\), \(M=N-1\), then uses \eqref{eq:split} for each
\(U_m\) with \(0\le m\le M\). A row is indexed by a subset.
For residue \(r\) list the pairs in \(a+b\le M\) with
\(2a+b\equiv r\pmod3\). Let \(K\) be the largest of the three list
lengths. Each row occupies \(K\) cells, with any padding zero. For each
destination residue and last label, store the pairs of predecessor and
destination slots in advance. If the destination is \((b,c)\) in
class \(r=\sum_{u\in S}u\bmod3\), label \(v\) uses the slot for
\((v-b-c,b)\) in the predecessor's own class \(r-v\bmod3\). At \(N=26\) the full triangle contains
\(\binom{27}{2}=351\) pairs, distributed \(117,117,117\).

For zero-based bit positions \(d_1<\cdots<d_h\) (label \(d_i+1\)),
the combinatorial number system gives the zero-based rank in
colexicographic order among subsets of size \(h\) as
\(\sum_{i=1}^h\binom{d_i}{i}\)~\cite{colex}, taking a binomial coefficient
as zero when the lower index exceeds the upper one. Grouping subsets by
their largest bit gives this formula recursively: \(\binom{d_h}{h}\)
subsets have a smaller largest bit. For deletion of \(d_j\) the rank becomes
\(\sum_{i<j}\binom{d_i}{i}+\sum_{i>j}\binom{d_i}{i-1}\).
One prefix and one suffix array therefore supply all deletion ranks in
\(O(h)\) arithmetic operations per row. The algorithm reads from layer
\(h-1\) and assigns each destination row of layer \(h\) to one OpenMP
worker. Row writes and the corresponding \(F\) writes are disjoint.
After the layer finishes only its successor needs the detailed weights.
The \(F\) array retains all \(2^M\) row totals for the final sums.
Its address is the local rank plus \(\sum_{i=0}^{h-1}\binom{M}{i}\).

Each ordering on \(h\) labels produces at most one path, so every cell,
row total, and nonnegative partial cell sum is bounded by \(h!\).
Since \(20!<2^{64}<21!\), the producer stores \(W\) cells in
\texttt{uint64\_t} through \(h=20\). At \(h=21\) it widens inputs
before addition, and from then on stores unsigned 128-bit cells.
All \(F\) values, products and final accumulators use 128 bits.
A complementary product counts a subclass of permutations, as does
each partial sum in \eqref{eq:split}, including the contribution obtained
by doubling a partial sum over unpaired representatives. The bound
\(n!\le28!<2^{98}\) covers the full implementation domain.
The type \texttt{unsigned \_\_int128} is a compiler extension on supported
targets~\cite{gcc}.

Writing \(s_h=8\) for \(h\le20\) and \(s_h=16\) otherwise, the principal
arrays at a layer transition use at most
\[
 16\,2^M+K\max_{0\le h<M}
 \left\{\binom{M}{h}s_h+\binom{M}{h+1}s_{h+1}\right\}
 \quad\text{bytes}\qquad(M\ge1).
\]
The triangle contains \(\binom{M+2}{2}\) pairs before slicing.
For \(M=25\) the array bound is \(10\,271\,832\,512\) bytes. A simple upper bound
on cell reads for the recurrence is \(K M2^{M-1}\), since each label belongs
to half the subsets. The precomputed edge lists reduce the actual number.
The convolution for one \(m>0\) requires only \(2^{m-1}\) products.

\section{A recount over alphabets with repetitions}\label{sec:ie}

Fix \(m\), and allow each label in \(T\subseteq U_m\) to be used any
number of times. Let \(L_j(T)\) count admissible words of length \(j\)
on this alphabet, including \(L_0(T)=1\). These counts require only the
length and the terminal pair. Initialize a single unit at \((0,0)\).
For each successive length, start a fresh zero row on the full triangle
and add the predecessor weight along \((a,b)\to(b,v-a-b)\) for
each \(v\in T\) satisfying \(v\ge a+b\).
The sum of the new row is \(L_j(T)\). The checker uses both coordinates
of the full triangle with its own numbering, without the residue partition
or the producer's subset ranks.

\begin{theorem}\label{thm:ie}
For every \(m\ge0\),
\begin{equation}
 a(m+1)=\sum_{T\subseteq U_m}(-1)^{m-|T|}
             \sum_{j=0}^{m} L_j(T)L_{m-j}(T).
 \label{eq:ie}
\end{equation}
\end{theorem}
\begin{proof}
The inner sum counts ordered pairs of words with total length \(m\).
Fix such a pair and let \(D\) be the set of all labels it actually uses.
It contributes for exactly the alphabets \(D\subseteq T\subseteq U_m\).
If a label is missing, toggling that label pairs equal contributions
of opposite signs. If \(D=U_m\), the sole contribution has sign \(+1\).
There are exactly \(m\) positions, so a surviving pair uses each of the
\(m\) labels exactly once. Theorem~\ref{thm:split} identifies it with one
feasible permutation. For \(m=0\) the two empty words give one.
\end{proof}

This is inclusion--exclusion on the Boolean lattice~\cite{rota}.
A related use of inclusion--exclusion counts Hamiltonian paths by
enforcing coverage in a walk count~\cite{karpie,bax}. In \eqref{eq:ie}
the objects to cover are the transition labels across two words. The
vertices are terminal pairs and may be visited repeatedly.

In a
job of maximum \(N\) each alphabet \(T\subseteq U_M\) supplies lengths
through \(M=N-1\) and contributes to each prefix \(m\) with
\(T\subseteq U_m\), equivalently \(\max(T)\le m\le M\), using
\(\max(\varnothing)=0\).
The C++ checker evaluates \eqref{eq:ie} in two arithmetic channels, one modulo each of
\(P=2^{61}-1\) and \(Q=2^{61}-3\). For either modulus, adding two
reduced residues fits in 64 unsigned bits and their product fits in
128. The two moduli are odd and differ by two, hence are coprime.
For reduced residues \(u\bmod P\) and \(v\bmod Q\), a separate Python
program returns
\[
 t=u+P\bigl((v-u)P^{-1}\bmod Q\bigr),\qquad 0\le t<PQ.
\]
It checks completed alphabet coverage and refuses \(t>n!\).
As \(28!<2^{98}<PQ\), the complete correct residues determine the exact
count over the entire supported range. For every reported \(n\), \(a(n)<Q<P\), so each final
residue already equals the displayed integer. Intermediate word counts
and alternating sums are computed in modular arithmetic. The two channels
use the same transitions and arithmetic routines. Agreement with the
separate weighted subset enumeration tests the count beyond what can
be learned from reconstructing the two residues alone.

\section{Computational result and checks}\label{sec:results}

\begin{theorem}[Computational result]\label{thm:result}
For \(17\le n\le26\), the numbers of permutations satisfying
\eqref{eq:def} are those in Table~\ref{tab:terms}.
\end{theorem}
\begin{proof}[Computational verification]
The weighted recurrence \eqref{eq:pull} and complementary sum \eqref{eq:split}
were evaluated for complete prefixes ending at 24 and 26. Separate runs
of \eqref{eq:ie} covered all \(2^{23}\) and \(2^{25}\) alphabets,
respectively. For each maximum, the complete reconstructed list agreed
entry by entry with the producer and reproduced the sixteen historical terms.
The second job includes \(a(25)\). The ten entries displayed here are the
corresponding rows of the resulting contiguous table.
\end{proof}

\begin{table}[htbp]\centering
\caption{Exact extension of the sixteen-term baseline for A180206.
Both blocks are indexed by the permutation length.}\label{tab:terms}
\begin{tabular}{rr@{\qquad}rr}\toprule
\(n\)&\(a(n)\)&\(n\)&\(a(n)\)\\\midrule
17 & 4379508168    & 22 & 170531561222304\\
18 & 33039187452   & 23 & 1606250624929548\\
19 & 261553618840  & 24 & 15740000714297474\\
20 & 2171808605286 & 25 & 160065864659464904\\
21 & 18838224530962& 26 & 1688253382074096582\\\bottomrule
\end{tabular}
\end{table}

\paragraph{Checks against the definition.}
Proposition~\ref{prop:direct} was evaluated on every permutation for
each \(1\le n\le8\), giving \(1,2,4,8,18,56,206,866\).
Every ordering of each of the 128 subsets of \(U_7\) was also examined,
checking the coefficients individually. An unpacked Python subset recurrence
with arbitrary-precision weights reproduced \(a(1),\ldots,a(16)\).
A separate C implementation of the difference criterion tested every
permutation through \(n=11\). Enumeration of all feasible starting pairs
for every permutation through 9 found no permutation with two realizations.
Further locally executed tests compared all \(2^{15}\) tail coefficients
on subsets of \(U_{15}\) with a separate implementation. A transfer computation
starting with every pair \((x_0,x_1)\) of sum at most \(n-1\) also
reproduced the first sixteen terms. These were the limits of the
definition and transfer checks available at consolidation.

\paragraph{A subsequent forward count.}
During manuscript review a separately written C program was run locally
for each \(1\le n\le21\). It begins with one empty history at every
nonnegative pair \((a,b)\) with \(a+b\le n-1\). Its subset records
which window sums in \(\{0,\ldots,n-1\}\), including zero, have been
used. From \((a,b)\) it appends any unused label \(v\ge a+b\) and
moves to \((b,v-a-b)\). Summing the terminal weights for the full
subset counts pairs of a permutation and its realization. Every realization
appears, since its first two entries sum to at most its first window.
A zero window forces three zeros and hence makes the realization unique.
This enumeration proceeds from left to right without splitting at zero.
All 21 totals agreed, including the first five new terms. Cell additions
were checked for overflow. Through \(n=20\) the final sum fits in 64 bits
by the factorial bound. At \(n=21\) a second run also checked every
addition in that final sum. A separate review program compared explicit
starting-pair enumeration with the difference criterion on every permutation
through \(n=11\), with no disagreement or multiple realization.

For \(22\le n\le26\) the complete numerical comparison reported here
is between the weighted subset and alphabet enumerations. They share the
zero decomposition and local transition. CRT reconstructs the outputs
of the alphabet enumeration.

\paragraph{Checks on the implementation.}
The release regression suite (101 invocations) compared one and eight requested workers against
the historical prefix through 16 and exercised interrupted and resumed
computations. A separate numerical fixture checked 768 synthetic rows
across cardinalities 20, 21 and 22. A forced accumulation of \(21!\)
exercised a value above the 64-bit limit. Clang and address/undefined-behaviour
sanitizer matrices were also exercised on the smaller test domain through 8.
Leak detection was disabled for the OpenMP address-sanitized tests.
The thread-sanitized tests used Archer, excluded noninstrumented modules,
and omitted the file-size-limit cases. A deliberate race and a race-free
control distinguished the two outcomes. A compatible single-worker
Valgrind build was checked at 8, and clang-tidy and Cppcheck diagnostics
were examined. These are finite software checks. The numerical evidence
for Table~\ref{tab:terms} comes from the two completed enumerations.

\paragraph{Execution and retained evidence.}
The computations used an AMD Ryzen 9 7940HS (8 cores, 16 logical CPUs)
under Linux Mint 22.3 (kernel \texttt{6.17.0-42-generic}), with eight
OpenMP workers requested. The GCC and
G++ 13.3.0 builds used C17 and C++17 language modes, respectively, with
\texttt{-O3 -march=native -fopenmp}. The C build also used
\texttt{-mtune=native -flto -funroll-loops -fomit-frame-pointer}.
The producer's startup test requires the array allocation estimate
plus \(8.5\,\mathrm{GiB}\) of available memory under Linux, even for small
jobs. Python 3.12.3 handled reconstruction. Table~\ref{tab:cost} gives the
measured costs. For example, at maximum 26 the checker performed
\(2M2^{M-1}\sum_{v=1}^{M}\binom{v+2}{2}
=2\,747\,269\,120\,000\) modular edge additions, for \(M=25\):
each alphabet containing \(v\) traverses its triangle of edges at
each length in both channels.

\begin{table}[htbp]\centering
\caption{Measured costs of complete prefix jobs. Times are rounded to whole
seconds. DP and IE are the weighted producer and the alphabet checker.
Peak resident set sizes refer to each individual numerical child process.
The group wall includes verification and consolidation. CPU sums user and
system time of the two numerical processes.}\label{tab:cost}
\small
\begin{tabular}{rrrrrrr}\toprule
 &\multicolumn{3}{c}{Wall time (s)}&CPU (s)&\multicolumn{2}{c}{Peak RSS (KiB)}\\
\cmidrule(lr){2-4}\cmidrule(lr){6-7}
\(N\)&DP&IE&Group&DP + IE&DP&IE\\\midrule
24&10&45&55&286&2204952&20176\\
26&46&243&289&1517&9855552&20212\\\bottomrule
\end{tabular}
\end{table}

Together the two jobs took about 344 seconds of group wall time,
excluding operator pauses and subsequent audit repetitions. Both native
programs checked their respective completed states before consolidation.
The producer checked committed spans of subset totals against their
recorded digests and bounds and compared the surviving last-layer row
totals with the corresponding subset totals. The checker verified its
committed state and coverage. Those four read-only checks were repeated
after the campaign, together
with one further CRT reconstruction at each maximum. The stored residues
and controls still permit reconstruction and comparison of every row.
The two full subset-count tables and the two surviving last-layer files
were subsequently discarded without an archive. The original producer
state can therefore no longer undergo its complete on-disk audit. The
retained result tables, program sources and validation records document
the computations described in the project reference~\cite{project}.

\section{What transfers to other window lengths?}\label{sec:widths}

For an unweighted window of length \(w\ge2\), its zero forces \(w\)
consecutive zeros. The outward update appends the next label minus the
sum of the preceding \(w-1\) entries. Thus A180205 uses a scalar and
A180207 uses a triple~\cite{hardintwo,hardinfour,two,four}. The decomposition into
two words applies to all three widths. The admissible words themselves
change with \(w\). At length 5 the three counts, in width order, are
\(44,18,16\). Even strict inequalities between counts at a fixed length
need not express containment.

For example, \(p=(0,1,2,4,3,6,5)\) is feasible at width three:
its right-hand appended entries are \((1,1,2,0,4,1)\). At widths two
and four the right updates first fail at the final label 5 and at label 3,
respectively, each requiring an entry of \(-1\).
Conversely, take \(p=(1,0,2,3,5,6,4)\). The left tail appends 1 at all
three widths. On the right, widths two and four append \((2,1,4,2,2)\)
and \((2,1,2,1,0)\). Width three appends \(2,1,2,3\), after which the
last label 4 would require \(4-2-3=-1\). Thus at the same length 7 the
width-three class is incomparable by inclusion with each neighbouring
class. The neighbouring counts cannot be substituted as counts of the
same objects.

\section{Conclusion}

The enumeration extends Hardin's sixteen-term baseline to a contiguous
table through \(a(26)=1688253382074096582\). Its central reduction is
that a zero sum determines the realization in both directions. What
remains to count on each side is an ordering of specified labels, with
two entries of memory and a computable residue class for their terminal
pair. Dropping distinctness within each tail and enforcing it globally by
inclusion--exclusion gives the same ten additional integers. The first
length beyond this computation is 27.

\section*{Data availability}
The author retains the indexed counts, original producer result lists,
modular outputs and controls, build identities, comparison receipts and
reconstruction code. The reconstruction program reads the exact checker
executable as data to verify its identity. It requires no execution of
that native binary. The stored records support comparison of both original
result lists with the integers reconstructed from the committed residues.
For the subsequent checks, the author retains both C sources, the exact
one-line change that checks the terminal sum, and the local execution
records. Both programs and the modified version were compiled with
GCC 13.3.0 using \texttt{-std=c17 -O2}.
The discarded subset tables and layers have no archive. Recomputing them
would generate new states. A public release of this project
has yet to be verified. The GitHub address in the References designates
the intended repository. The Zenodo DOI placeholder will be replaced
after the first Zenodo publication when the author supplies the verified
official locator.

\section*{AI use disclosure}
OpenAI Codex assisted with mathematical analysis, implementation and
drafting. Code review also drew on reports from OpenAI and from tools
identified in the project records as Claude, GLM, MiniMax, Qwen and
Meta AI (self-reported Muse family). Additional comments were received through Gemini and
Grok. Some deployed model versions could not be established.

\noindent\textit{Mathematical analysis.}
Codex and a NotebookLM consultation helped compare the adjacent-window
methods and formulate checks of the pair recurrence. The congruence and
the examples distinguishing label sets were verified from the transition.
The historical prefix was checked against Hardin's entry.

\noindent\textit{Algorithm and code.}
Codex assisted with the weighted producer, the modular recount and
their tests. Code-review reports drew attention to resource checks on
restart, checkpoint coverage and the integer-width transition. Numerical
claims rest on the local result tables, stored residues, and executed
comparisons. They do not rely on conversational assertions or reviewer
agreement.

\noindent\textit{Manuscript.}
Codex assisted with exposition, literature checks and the comparison of
proofs, table entries and computational claims with the source programs
and original execution records. External manuscript reviews identified
the methodological references and suggested clarifications to the evidence
and attribution. The two review programs used for the subsequent checks
were supplied through Claude. Their numerical results were checked by
local execution. Carlo Corti takes full scientific
responsibility for the mathematical statements, algorithm, implementation,
numerical results, validation boundaries and conclusions.

\clearpage
\begin{thebibliography}{99}
\raggedright

\bibitem{hardin}
R.~H. Hardin, Number of permutations of \(0..(n-1)\) representable as
consecutive sums of 3 adjacent elements of a sequence of \(n+2\)
nonnegative integers, A180206, \emph{The On-Line Encyclopedia of Integer
Sequences}, 15 August 2010. \url{https://oeis.org/A180206}.
Indexed table consulted 18 November 2026,
\url{https://oeis.org/A180206/b180206.txt}.
Official data mirror, revision 10, 8 November 2024,
\url{https://github.com/oeis/oeisdata/blob/main/seq/A180/A180206.seq}.

\bibitem{hardintwo}
R.~H. Hardin, OEIS entry A180205 (adjacent pairs), 15 August 2010.
\url{https://oeis.org/A180205}.

\bibitem{hardinfour}
R.~H. Hardin, OEIS entry A180207 (windows of length four), 15 August 2010.
Revision 15, 15 November 2026, credits C. Corti for terms 17--22.
\url{https://oeis.org/A180207/internal}.

\bibitem{two}
C. Corti, \emph{A180205: exact enumeration through \(n=28\)},
Zenodo release series (adjacent-sum permutations), 2026.
\url{https://doi.org/10.5281/zenodo.24311714}.

\bibitem{four}
C. Corti, \emph{A180207: exact enumeration through \(n=22\)},
Zenodo release series (four-window permutations), 2026.
\url{https://doi.org/10.5281/zenodo.24194032}.

\bibitem{forward}
R.~H. Hardin, Number of arrays of length \(n\) that are sums of 3
consecutive elements of length \(n+2\) permutations of \(0..n+1\), A229561, \emph{The On-Line Encyclopedia of Integer Sequences},
26 September 2013. \url{https://oeis.org/A229561}.

\bibitem{bellman}
R. Bellman, Dynamic programming treatment of the travelling salesman
problem, \emph{J. ACM} \textbf{9} (1962), 61--63.

\bibitem{held}
M. Held and R.~M. Karp, A dynamic programming approach to sequencing
problems, \emph{J. Soc. Indust. Appl. Math.} \textbf{10} (1962), 196--210.
\url{https://doi.org/10.1137/0110015}.

\bibitem{karpie}
R.~M. Karp, Dynamic programming meets the principle of inclusion and
exclusion, \emph{Oper. Res. Lett.} \textbf{1} (1982), 49--51.

\bibitem{bax}
E.~T. Bax, Inclusion and exclusion algorithm for the Hamiltonian path
problem, \emph{Inform. Process. Lett.} \textbf{47} (1993), 203--207.

\bibitem{colex}
V. Kruchinin, Y. Shablya, D. Kruchinin and V. Rulevskiy,
Unranking Small Combinations of a Large Set in Co-Lexicographic Order,
\emph{Algorithms} \textbf{15} (2022), no.~2, article 36.
\url{https://doi.org/10.3390/a15020036}.

\bibitem{gcc}
GNU Project, \emph{Using the GNU Compiler Collection}, ``128-bit Integers,''
consulted November 2026.
\url{https://gcc.gnu.org/onlinedocs/gcc/_005f_005fint128.html}.

\bibitem{rota}
G.-C. Rota, On the foundations of combinatorial theory I. Theory of
M\"obius functions, \emph{Z. Wahrscheinlichkeitstheorie verw. Gebiete}
\textbf{2} (1964), 340--368.
\url{https://doi.org/10.1007/BF00531932}.

\bibitem{project}
C. Corti, \emph{Counting Permutations of Consecutive Triple Sums:
OEIS A180206 through \(n=26\)}, GitHub repository and Zenodo archive, 2026.
\url{https://github.com/carcorti/A180206}.
\url{https://doi.org/10.5281/zenodo.xxxxxxxx}.
Intended release locators. The DOI placeholder awaits the first publication
on Zenodo and the author's verified official locator.
\end{thebibliography}
\end{document}
